\section{带边界流形}\label{sec:manifolds-with-boundary}

再次考虑\textbf{闭半空间}
\[
  \mathbb{H}^n=\{(x^1,\ldots,x^n)\in\mathbb{R}^n\mid x^n\geq 0\}
\]
及由 $\mathbb{R}^n$ 的通常拓扑诱导的拓扑. 回顾一下, 定义在开集 $U\subset\mathbb{H}^n$ 上的映射 $f:U\to\mathbb{R}^m$ 称为\textbf{可微}, 若它是定义在 $\mathbb{R}^n$ 中包含 $U$ 的开子集上的某个可微映射 $\overline f$ 到 $U$ 的限制 (参见第~\RGref{subsec:notes-differentiability}~节). 此时定义导数 $(\mathrm{d}f)_p$ 为 $(\mathrm{d}\overline f)_p$. 注意, 这个导数与所用的延拓无关, 因为任何两个延拓都必须在 $U$ 上一致.

\begin{RiemannDefinition}\label{def:smooth-manifold-with-boundary}
\textbf{带边界的光滑 $n$ 维流形}是一个 $n$ 维带边界拓扑流形 $M$ 和一族参数化 $\varphi_\alpha:U_\alpha\subset\mathbb{H}^n\to M$ (即从 $\mathbb{H}^n$ 的开集 $U_\alpha$ 到 $M$ 的开集的同胚), 使得:
\begin{enumerate}[label=(\roman*),leftmargin=*]
\item 坐标邻域覆盖 $M$, 即 $\bigcup_\alpha\varphi_\alpha(U_\alpha)=M$;
\item 对每一对满足
\[
  W:=\varphi_\alpha(U_\alpha)\cap\varphi_\beta(U_\beta)\ne\varnothing,
\]
的指标 $\alpha,\beta$, 重叠映射
\begin{gather*}
  \varphi_\beta^{-1}\circ\varphi_\alpha:\varphi_\alpha^{-1}(W)\to\varphi_\beta^{-1}(W)\\
  \varphi_\alpha^{-1}\circ\varphi_\beta:\varphi_\beta^{-1}(W)\to\varphi_\alpha^{-1}(W)
\end{gather*}
都是光滑的;
\item 族 $\mathcal{A}=\{(U_\alpha,\varphi_\alpha)\}$ 关于 (i) 和 (ii) 是极大的, 即若 $\varphi_0:U_0\to M$ 是一个参数化, 使得对 $\mathcal{A}$ 中所有 $\varphi$, $\varphi_0\circ\varphi^{-1}$ 和 $\varphi^{-1}\circ\varphi_0$ 都是 $C^\infty$ 的, 则 $\varphi_0$ 属于 $\mathcal{A}$.
\end{enumerate}
\end{RiemannDefinition}

回顾一下, 若 $M$ 中一点属于 $\partial\mathbb{H}^n$ 在某个参数化下的像 (即存在参数化 $\varphi:U\subset\mathbb{H}^n\to M$, 对某个 $(x^1,\ldots,x^{n-1})\in\mathbb{R}^{n-1}$ 有 $\varphi(x^1,\ldots,x^{n-1},0)=p$), 则称它为\textbf{边界点}. 所有这些点所成的集合 $\partial M$ 称为 $M$ 的\textbf{边界}. 注意, 可微流形是带边界可微流形的特殊情形, 对此有 $\partial M=\varnothing$.

\begin{RiemannProposition}\label{prop:smooth-boundary-manifold}
带边界的光滑 $n$ 维流形的边界是一个 $(n-1)$ 维可微流形.
\end{RiemannProposition}

\begin{RiemannProof}
设 $p$ 是 $M$ (一个带边界 $n$ 维流形) 的边界点, 在 $p$ 附近选取参数化 $\varphi_\alpha:U_\alpha\subset\mathbb{H}^n\to M$. 令 $V_\alpha:=\varphi_\alpha(U_\alpha)$, 我们断言 $\varphi_\alpha(\partial U_\alpha)=\partial V_\alpha$, 其中 $\partial U_\alpha=U_\alpha\cap\partial\mathbb{H}^n$ 且 $\partial V_\alpha=V_\alpha\cap\partial M$. 根据边界的定义, 已知 $\varphi_\alpha(\partial U_\alpha)\subset\partial V_\alpha$, 所以只需证明 $\partial V_\alpha\subset\varphi_\alpha(\partial U_\alpha)$. 设 $q\in\partial V_\alpha$, 考虑 $q$ 附近的参数化 $\varphi_\beta:U_\beta\to V_\alpha$, 它将 $\mathbb{H}^n$ 的一个开子集映到 $M$ 的一个开子集, 且使 $q\in\varphi_\beta(\partial U_\beta)$. 如果能证明 $\varphi_\beta(\partial U_\beta)\subset\varphi_\alpha(\partial U_\alpha)$, 结论便成立. 为此, 证明 $(\varphi_\alpha^{-1}\circ\varphi_\beta)(\partial U_\beta)\subset\partial U_\alpha$. 事实上, 假设这个映射 $\varphi_\alpha^{-1}\circ\varphi_\beta$ 将一点 $x\in\partial U_\beta$ 映到 $U_\alpha$ 的一个 (在 $\mathbb{R}^n$ 中的) 内点. 由于这个映射是微分同胚, $x$ 就是 $U_\beta$ 的一个 (在 $\mathbb{R}^n$ 中的) 内点. 这当然与假设 $x\in\partial U_\beta$ 矛盾. 因此, $(\varphi_\alpha^{-1}\circ\varphi_\beta)(\partial U_\beta)\subset\partial U_\alpha$, 从而 $\varphi_\beta(\partial U_\beta)\subset\varphi_\alpha(\partial U_\alpha)$.

限制 $\varphi_\alpha$ 就得到从 $\partial U_\alpha$ 到 $\partial V_\alpha$ 的微分同胚, 其中把 $\partial U_\alpha$ 与 $\mathbb{R}^{n-1}$ 的一个开子集等同. 这样就得到了 $\partial M$ 中 $p$ 附近的一个参数化, 容易看出这些参数化在 $\partial M$ 上定义了可微结构.
\end{RiemannProof}

\begin{RiemannRemark}\label{rem:boundary-point-independent-of-chart}
由上面的证明可知, 边界点的定义与所选参数化无关. 也就是说, 若存在 $p$ 附近的一个参数化, 使得 $p$ 为 $\partial\mathbb{H}^n$ 中一点的像, 那么 $p$ 附近的任何参数化都将 $\mathbb{H}^n$ 的一个边界点映到 $p$.
\end{RiemannRemark}

可定向性的定义很容易推广到带边界流形. 此时有下面的结果.

\begin{RiemannProposition}\label{prop:orientability-of-boundary}
设 $M$ 是可定向带边界流形. 则 $\partial M$ 也可定向.
\end{RiemannProposition}

\begin{RiemannProof}
若 $M$ 可定向, 可以在 $M$ 上选取一个图册 $\{(U_\alpha,\varphi_\alpha)\}$, 使得所有重叠映射的导数的行列式均为正. 利用这个图册, 可以按 Proposition~\ref{prop:smooth-boundary-manifold} 的证明中所述的方法, 得到 $\partial M$ 的图册 $\{(\partial U_\alpha,\widetilde\varphi_\alpha)\}$. 对任意重叠映射
\[
  (\varphi_\beta^{-1}\circ\varphi_\alpha)(x^1,\ldots,x^n)
  =(y^1(x^1,\ldots,x^n),\ldots,y^n(x^1,\ldots,x^n))
\]
有
\begin{align*}
  &(\varphi_\beta^{-1}\circ\varphi_\alpha)(x^1,\ldots,x^{n-1},0)\\
  &\quad=(y^1(x^1,\ldots,x^{n-1},0),\ldots,y^{n-1}(x^1,\ldots,x^{n-1},0),0)
\end{align*}
且
\begin{align*}
  &(\widetilde\varphi_\beta^{-1}\circ\widetilde\varphi_\alpha)(x^1,\ldots,x^{n-1})\\
  &\quad=(y^1(x^1,\ldots,x^{n-1},0),\ldots,y^{n-1}(x^1,\ldots,x^{n-1},0)).
\end{align*}
于是, 将 $(x^1,\ldots,x^{n-1},0)$ 记为 $(\widetilde x,0)$, 有
\[
  \bigl(\mathrm{d}(\varphi_\beta^{-1}\circ\varphi_\alpha)\bigr)_{(\widetilde x,0)}
   =\left(\begin{array}{c:c}
    \bigl(\mathrm{d}(\widetilde\varphi_\beta^{-1}\circ\widetilde\varphi_\alpha)\bigr)_{\widetilde x}&*\\
    \hdashline
    0&\frac{\partial y^n}{\partial x^n}(\widetilde x,0)
  \end{array}\right)
\]
从而
\[
  \det\bigl(\mathrm{d}(\varphi_\beta^{-1}\circ\varphi_\alpha)\bigr)_{(\widetilde x,0)}
  =\frac{\partial y^n}{\partial x^n}(\widetilde x,0)
    \det\bigl(\mathrm{d}(\widetilde\varphi_\beta^{-1}\circ\widetilde\varphi_\alpha)\bigr)_{\widetilde x}.
\]
然而, 固定 $x^1,\ldots,x^{n-1}$ 后, 当 $x^n$ 为正时 $y^n$ 为正, 当 $x^n=0$ 时 $y^n$ 为零. 因此 $\frac{\partial y^n}{\partial x^n}(\widetilde x,0)>0$, 所以
\[
  \det\bigl(\mathrm{d}(\widetilde\varphi_\beta^{-1}\circ\widetilde\varphi_\alpha)\bigr)_{\widetilde x}>0.
\]
\end{RiemannProof}

因此, 在带边界流形 $M$ 上选取一个定向, 会在边界 $\partial M$ 上诱导出一个定向. 一种方便的选择称为\textbf{诱导定向}, 可以按如下方式得到. 对 $p\in\partial M$, 切空间 $T_p(\partial M)$ 是 $T_pM$ 的余维为 $1$ 的子空间. 如上所见, 考虑 $p$ 附近的坐标系 $x:W\to\mathbb{R}^n$, 有 $x^n(p)=0$, 且 $(x^1,\ldots,x^{n-1})$ 是 $\partial M$ 中 $p$ 附近的坐标系. 令 $n_p:=-\bigl(\frac{\partial}{\partial x^n}\bigr)_p$ (称为 $p$ 处的\textbf{向外向量}). 若 $T_pM$ 的有序基 $\{n_p,\beta\}$ 为正, 就给 $T_p(\partial M)$ 的有序基 $\beta$ 赋予正号, 否则赋予负号. 由此定义 $\partial M$ 上的诱导定向. 注意, 由于 $\frac{\partial y^n}{\partial x^n}(\varphi^{-1}(p))>0$ (采用上面的记号), $n_p$ 的最后一个分量的符号与坐标系的选取无关. 一般而言, 诱导定向\textbf{不是}从 $M$ 的坐标图中直接去掉最后一个坐标所得的定向 (实际上, 它是该定向的 $(-1)^n$ 倍).

\begin{RiemannExercise}\label{ex:manifolds-with-boundary}
\begin{enumerate}[label=(\arabic*),ref=\arabic*,leftmargin=*]
\item\label{ex:product-with-corners} 举例说明两个带边界流形的乘积不总是带边界流形.
\item\label{ex:product-with-boundaryless-manifold} 设 $M$ 是无边界流形, $N$ 是带边界流形. 证明乘积 $M\times N$ 是带边界流形. $\partial(M\times N)$ 是什么?
\item\label{ex:diffeomorphism-preserves-boundary} 证明两个带边界流形 $M$ 和 $N$ 之间的微分同胚将边界 $\partial M$ 微分同胚地映到 $\partial N$.
\end{enumerate}
\end{RiemannExercise}
